\paragraph{Part 2: Finding a $g,h$ pair within Candidates}
First, we note that any output by Algorithm 2 must be correct.
\begin{ntclaim}[Author Claim 39]
Any value Algorithm 2 outputs must be an element of $\Opt(f)$.
\end{ntclaim}
\begin{proof}
Any output $\{g,O,h\}$ of Algorithm 2 must satisfy $f=gOh$ and $L(f)=L(g)+L(h)$ which implies $\{g,O,h\}\in\Opt(f)$ by Lemma 21.
\end{proof}
Finally, we show that Algorithm 2 must output a value.
\begin{ntclaim}[Author Claim 40]
Algorithm 2 must output a value (on input $f$).
\end{ntclaim}
\begin{proof}
Since $L(f)\geq 2$, we have that $\Opt(f)$ is non-empty. Let $\{g,O,h\}\in\Opt(f)$.
Claim 38 implies that $\{g,h\}\subseteq\Candidates$. On the other hand, Lemma 21 implies that $L(f)=L(g)+L(h)$ and $f=gOh$. Thus, it is clear that part 2 will either output $\{g,O,h\}$ or output a value before that.
\end{proof}
\subsection*{5.2 Running Time of Algorithm 2}
Fix some function $f$ with $n$-inputs satisfying $L(f)\geq 2$. Let $N=2^n$. We break the running time analysis into the two pieces of the algorithm.
\paragraph{Part 1}
It is easy to see that the run time of part 1 can be bounded by
\[O\left(N+\sum_{i\in[N]}N\cdot|\PC{i}|\right)\]
where $|\PC{i}|$ indicates the size of $\PC{i}$ after Algorithm 2 is finished adding elements to it.
Moreover, we can bound the quantity $|\PC{i}|$ as follows. Claim 37 implies that every partial function in $\PC{i}$ must be consistent with some total function $g$ on $n$-inputs satisfying $L(\Select[f,g])\leq L(f)+n+2$.
On the other hand, Claim 36 implies that any single (total) function can agree with at most partial function in $\PC{i}$. Hence, we have that
\[|\PC{i}|\leq|\{g:L(\Select[f,g])\leq L(f)+n+2\}|\]
and Lemma 24 implies that
\[|\{g:L(\Select[f,g])\leq L(f)+n+2\}|\leq O(|\Canon{n+1}(f)|\cdot N^2\log N).\]
Thus, we have that part 1 runs in time at most $O(|\Canon{n+1}(f)|\cdot N^4\log N)$. Moreover, part 1 only makes oracle calls of length at most $2N$ (to calculate $L(\Select[f,g_{\gamma_b}])$).
\paragraph{Part 2}
It is easy to see that this part runs in time $O(N\cdot|\Candidates|^2)$. Hence, since $\Candidates=\PC{N}$, the analysis in part 1 above implies that
\[|\Candidates|\leq O(|\Canon{n+1}(f)|\cdot N^2\log N).\]
Thus, part 2 runs in time at most
\[O(|\Canon{n+1}(f)|^2\cdot N^5\log^2N).\]
Moreover, part 2 only makes oracle calls of length $N$.
\paragraph{In total}
Putting it all together, we have that Algorithm 2 runs in time at most
\[O(|\Canon{n+1}(f)|^2\cdot N^5\log^2N)\]
and only makes oracle queries of length $2N$.
\end{proof}
